% ==============================================================================
% SECCIÓN 1: OBJETIVO Y FUNDAMENTO TEÓRICO
% ==============================================================================

\section{Objetivos del Laboratorio}

\subsection{Objetivo General}
Definir formalmente la visión del sistema, el mapa integral de \textit{stakeholders}, los objetivos estratégicos y el alcance inicial del proyecto \textbf{“Sistema Integral de Gestión de Capacitaciones (SIGC)”}, empleando técnicas rigurosas de \textit{Agile Inception}, modelado de negocio mediante \textit{Lean Canvas} y gestión visual de artefactos con la plataforma \textbf{Trello}.

\subsection{Objetivos Específicos}
\begin{itemize}
    \item Identificar, clasificar y priorizar a los actores clave (\textit{stakeholders}) del ecosistema formativo (entidades públicas como municipalidades, empresas privadas del sector turismo, capacitadores y participantes).
    \item Construir la declaración canónica de visión del producto utilizando el estándar de Geoffrey Moore, garantizando un posicionamiento claro y diferenciado frente a métodos manuales.
    \item Desarrollar el modelo de negocio y viabilidad funcional del producto en un \textit{Lean Canvas} integral de 9 bloques, diseñado en formato horizontal para maximizar su legibilidad analítica.
    \item Formular una matriz de riesgos tempranos con sus respectivas estrategias de mitigación técnica y operativa (enfocadas en conectividad, adopción de usuarios y autenticidad documental).
    \item Establecer la estructura de gestión ágil en \textbf{Trello} para el seguimiento de épicas, historias de usuario y actividades de desarrollo de la versión \textbf{SIGC\_V0.1}.
\end{itemize}

\section{Fundamento Teórico}

\subsection{Agile Inception y Descubrimiento de Producto}
La fase de \textit{Agile Inception} (o \textit{Inception Deck}), popularizada por Rasmusson y adaptada en el marco de \textit{Lean Inception} por Caroli, representa el punto de partida en los proyectos ágiles de ingeniería de software. Su propósito no es elaborar especificaciones extensas o cerradas, sino alinear a los desarrolladores, líderes de proyecto y partes interesadas en torno a cuatro preguntas fundamentales antes de escribir código:
\begin{enumerate}
    \item \textbf{¿Por qué estamos construyendo este sistema?} (Propósito, dolor real y justificación de negocio).
    \item \textbf{¿Para quién lo construimos?} (Segmentos de usuarios y beneficiarios directos).
    \item \textbf{¿Qué es y qué no es el producto?} (Fronteras y alcance inicial controlable o MVP).
    \item \textbf{¿Cómo mediremos el éxito y mitigaremos las incertidumbres?} (Métricas clave y riesgos tempranos).
\end{enumerate}

\subsection{Lean Canvas (Metodología de Ash Maurya)}
El \textit{Lean Canvas} es una adaptación orientada a software y startups del clásico \textit{Business Model Canvas} de Alexander Osterwalder. Creado por Ash Maurya, reemplaza componentes tradicionales de infraestructura y socios por bloques críticos de validación de riesgo:
\begin{itemize}
    \item \textbf{Problema y Alternativas Actuales:} Identifica las tres principales carencias operativas y cómo se resuelven de forma precaria hoy en día.
    \item \textbf{Segmentos de Clientes y \textit{Early Adopters}:} Define a los beneficiarios y usuarios que sufren con mayor urgencia el problema.
    \item \textbf{Propuesta de Valor Única (UVP):} Mensaje conciso que promete un beneficio transformador.
    \item \textbf{Solución:} Las características tecnológicas de software esenciales para abordar cada problema.
    \item \textbf{Canales, Flujos de Ingresos y Estructura de Costes:} Sostenibilidad operativa y económica del proyecto.
    \item \textbf{Métricas Clave y Ventaja Competitiva (\textit{Unfair Advantage}):} Indicadores de salud del sistema y barreras de entrada.
\end{itemize}

\subsection{Gestión de Stakeholders según Estándares de Ingeniería de Software}
De acuerdo con el estándar ISO/IEC/IEEE 29148 (\textit{Requirements Engineering}), un \textit{stakeholder} es cualquier individuo, grupo u organización que puede afectar, ser afectado o percibirse como afectado por una decisión, actividad o resultado del sistema. En el desarrollo de software moderno, una omisión en la identificación temprana de actores conduce a requerimientos incompletos, fallos de adopción y sobrecostes de reingeniería.

\subsection{Declaración de Visión según Geoffrey Moore}
La formulación de la visión del producto sigue la estructura canónica de Geoffrey Moore (\textit{Crossing the Chasm}):
\begin{quote}
    \textit{“Para [cliente/usuario objetivo] que [necesidad o dolor], el sistema [nombre] es un [categoría de producto] que [beneficio clave]. A diferencia de [alternativas o competidores actuales], nuestro producto [diferenciador principal]”.}
\end{quote}

\section{Herramientas Empleadas}
Para la ejecución integral del laboratorio se emplean:
\begin{itemize}
    \item \textbf{Trello:} Herramienta principal de gestión ágil basada en tableros Kanban, utilizada para estructurar el \textit{Product Backlog}, los \textit{Sprints} de requerimientos y el flujo de trabajo del equipo.
    \item \textbf{Miro / Notion:} Entornos colaborativos para la ideación interactiva y el registro de minutas técnicas.
    \item \textbf{\LaTeX{} (MiKTeX/TeX Live):} Sistema de composición tipográfica empleado para producir la documentación técnica bajo los más estrictos estándares académicos y normativos.
\end{itemize}
